\chapter{Implementation}

This chapter details the practical implementation of the research methodology, including data preprocessing pipelines, baseline model training procedures, and the AHFS-TA framework development.

\section{Development Environment}

\textbf{Hardware Configuration}:
\begin{itemize}
\item CPU: Intel Core i7 / AMD Ryzen 7 (8 cores)
\item RAM: 16 GB DDR4
\item GPU: NVIDIA GeForce GTX 1660 Ti / RTX 3060 (for deep learning)
\item Storage: 512 GB SSD
\end{itemize}

\textbf{Software Stack}:
\begin{itemize}
\item Operating System: Windows 10/11, Ubuntu 20.04 LTS
\item Python Version: 3.10.10
\item IDE: VS Code, Jupyter Notebook
\item Version Control: Git 2.40
\end{itemize}

\textbf{Python Libraries}:
\begin{verbatim}
pandas==2.0.3
numpy==1.24.3
scikit-learn==1.3.0
xgboost==2.0.0
torch==2.0.1
transformers==4.30.2
shap==0.42.1
matplotlib==3.7.2
seaborn==0.12.2
\end{verbatim}

\section{Data Preprocessing Pipeline}

\subsection{Data Loading}

\begin{verbatim}
import pandas as pd
import numpy as np

# Load dataset
data = pd.read_csv('data/educational_data.csv')
print(f"Total students: {len(data)}")
print(f"Total features: {data.shape[1]}")
print(f"Classes: {data['Target'].unique()}")
\end{verbatim}

\textbf{Output}: 4,424 students, 37 columns (36 features + 1 target), 3 classes.

\subsection{Missing Value Handling}

\begin{verbatim}
# Check missing values
missing = data.isnull().sum()
print(f"Features with missing values: {(missing > 0).sum()}")

# Imputation strategy
from sklearn.impute import SimpleImputer

# Numerical features: Median imputation
num_imputer = SimpleImputer(strategy='median')
num_features = data.select_dtypes(include=[np.number]).columns
data[num_features] = num_imputer.fit_transform(data[num_features])

# Categorical features: Mode imputation
cat_imputer = SimpleImputer(strategy='most_frequent')
cat_features = data.select_dtypes(include=['object']).columns
data[cat_features] = cat_imputer.fit_transform(data[cat_features])
\end{verbatim}

\subsection{Feature Encoding}

\begin{verbatim}
from sklearn.preprocessing import LabelEncoder, StandardScaler

# Encode target variable
le = LabelEncoder()
data['Target_Encoded'] = le.fit_transform(data['Target'])
# 0: Dropout, 1: Enrolled, 2: Graduate

# Encode categorical features
for col in cat_features:
    if col != 'Target':
        data[f'{col}_encoded'] = le.fit_transform(data[col])

# Drop original categorical columns
X = data.drop(['Target', 'Target_Encoded'] + list(cat_features), axis=1)
y = data['Target_Encoded']
\end{verbatim}

\subsection{Feature Scaling}

\begin{verbatim}
# Standardization for neural networks
scaler = StandardScaler()
X_scaled = scaler.fit_transform(X)

# Min-Max scaling for tree-based models (optional)
from sklearn.preprocessing import MinMaxScaler
mm_scaler = MinMaxScaler()
X_minmax = mm_scaler.fit_transform(X)
\end{verbatim}

\subsection{Train-Test Split}

\begin{verbatim}
from sklearn.model_selection import train_test_split

# Stratified 80/20 split
X_train, X_test, y_train, y_test = train_test_split(
    X_scaled, y, test_size=0.2, 
    stratify=y, random_state=42
)

print(f"Training set: {len(X_train)} students")
print(f"Test set: {len(X_test)} students")
print(f"Class distribution (train): {np.bincount(y_train)}")
print(f"Class distribution (test): {np.bincount(y_test)}")
\end{verbatim}

\textbf{Output}:
\begin{itemize}
\item Training: 3,539 students
\item Test: 885 students
\item Train distribution: [1,137 Dropout, 635 Enrolled, 1,767 Graduate]
\item Test distribution: [284 Dropout, 159 Enrolled, 442 Graduate]
\end{itemize}

\section{Feature Ranking Implementation}

\subsection{Information Gain and Gain Ratio}

\begin{verbatim}
from sklearn.feature_selection import mutual_info_classif

# Information Gain (Mutual Information)
ig_scores = mutual_info_classif(X_train, y_train, random_state=42)
ig_ranking = pd.DataFrame({
    'Feature': X.columns,
    'IG_Score': ig_scores
}).sort_values('IG_Score', ascending=False)

# Gain Ratio (IG normalized by feature entropy)
from scipy.stats import entropy

def calculate_gain_ratio(X, y):
    ig = mutual_info_classif(X, y, random_state=42)
    gr = []
    for i in range(X.shape[1]):
        feature_entropy = entropy(np.histogram(X[:, i], bins=10)[0] + 1e-10)
        gr.append(ig[i] / (feature_entropy + 1e-10))
    return np.array(gr)

gr_scores = calculate_gain_ratio(X_train, y_train)
\end{verbatim}

\subsection{Gini Importance}

\begin{verbatim}
from sklearn.ensemble import RandomForestClassifier

# Train Random Forest for Gini importances
rf = RandomForestClassifier(
    n_estimators=200, max_depth=15, 
    random_state=42, n_jobs=-1
)
rf.fit(X_train, y_train)

gini_scores = rf.feature_importances_
\end{verbatim}

\subsection{Chi-squared and F-statistic}

\begin{verbatim}
from sklearn.feature_selection import chi2, f_classif

# Chi-squared test (for non-negative features)
X_train_nonneg = X_train - X_train.min() + 1e-10
chi2_scores, chi2_pvalues = chi2(X_train_nonneg, y_train)

# ANOVA F-statistic
f_scores, f_pvalues = f_classif(X_train, y_train)
\end{verbatim}

\subsection{Unified Ranking Table}

\begin{verbatim}
# Combine all rankings
ranking_df = pd.DataFrame({
    'Feature': X.columns,
    'IG': ig_scores,
    'GR': gr_scores,
    'Gini': gini_scores,
    'Chi2': chi2_scores,
    'F': f_scores
})

# Rank each method
for method in ['IG', 'GR', 'Gini', 'Chi2', 'F']:
    ranking_df[f'{method}_Rank'] = ranking_df[method].rank(
        ascending=False, method='dense'
    )

# Average rank
ranking_df['Average_Rank'] = ranking_df[
    ['IG_Rank', 'GR_Rank', 'Gini_Rank', 'Chi2_Rank', 'F_Rank']
].mean(axis=1)

# Final ranking
ranking_df['Final_Rank'] = ranking_df['Average_Rank'].rank(method='dense')
ranking_df = ranking_df.sort_values('Final_Rank')

print(ranking_df.head(20))
\end{verbatim}

\section{Baseline Model Training}

\subsection{Decision Tree}

\begin{verbatim}
from sklearn.tree import DecisionTreeClassifier
from sklearn.model_selection import GridSearchCV

# Hyperparameter tuning
param_grid = {
    'max_depth': [5, 10, 15, 20],
    'min_samples_split': [10, 20, 30],
    'min_samples_leaf': [5, 10, 15]
}

dt = DecisionTreeClassifier(criterion='gini', random_state=42)
grid_search = GridSearchCV(
    dt, param_grid, cv=5, scoring='accuracy', n_jobs=-1
)
grid_search.fit(X_train, y_train)

# Best model
best_dt = grid_search.best_estimator_
print(f"Best params: {grid_search.best_params_}")

# Predictions
y_pred_dt = best_dt.predict(X_test)
\end{verbatim}

\subsection{Naive Bayes}

\begin{verbatim}
from sklearn.naive_bayes import GaussianNB

nb = GaussianNB()
nb.fit(X_train, y_train)
y_pred_nb = nb.predict(X_test)
\end{verbatim}

\subsection{Random Forest}

\begin{verbatim}
rf = RandomForestClassifier(
    n_estimators=200, max_depth=15, 
    max_features='sqrt', min_samples_split=10,
    random_state=42, n_jobs=-1
)
rf.fit(X_train, y_train)
y_pred_rf = rf.predict(X_test)
\end{verbatim}

\subsection{AdaBoost}

\begin{verbatim}
from sklearn.ensemble import AdaBoostClassifier

adaboost = AdaBoostClassifier(
    base_estimator=DecisionTreeClassifier(max_depth=1),
    n_estimators=100, learning_rate=0.5, random_state=42
)
adaboost.fit(X_train, y_train)
y_pred_ada = adaboost.predict(X_test)
\end{verbatim}

\subsection{XGBoost}

\begin{verbatim}
import xgboost as xgb

xgb_model = xgb.XGBClassifier(
    n_estimators=200, max_depth=6, learning_rate=0.1,
    subsample=0.8, colsample_bytree=0.8,
    reg_alpha=0.1, reg_lambda=1.0,
    random_state=42, n_jobs=-1
)
xgb_model.fit(X_train, y_train)
y_pred_xgb = xgb_model.predict(X_test)
\end{verbatim}

\subsection{Neural Network}

\begin{verbatim}
from sklearn.neural_network import MLPClassifier

nn = MLPClassifier(
    hidden_layer_sizes=(128, 64, 32),
    activation='relu', solver='adam',
    learning_rate_init=0.001, max_iter=100,
    early_stopping=True, validation_fraction=0.2,
    random_state=42
)
nn.fit(X_train, y_train)
y_pred_nn = nn.predict(X_test)
\end{verbatim}

\section{Model Evaluation}

\subsection{Performance Metrics Calculation}

\begin{verbatim}
from sklearn.metrics import (
    accuracy_score, precision_recall_fscore_support,
    confusion_matrix, roc_auc_score, classification_report
)

def evaluate_model(y_true, y_pred, y_prob, model_name):
    # Accuracy
    acc = accuracy_score(y_true, y_pred)
    
    # Precision, Recall, F1 (macro-averaged)
    precision, recall, f1, _ = precision_recall_fscore_support(
        y_true, y_pred, average='macro'
    )
    
    # AUC-ROC (one-vs-rest for multi-class)
    auc = roc_auc_score(y_true, y_prob, multi_class='ovr', average='macro')
    
    # Confusion matrix
    cm = confusion_matrix(y_true, y_pred)
    
    # Print results
    print(f"\n{model_name} Performance:")
    print(f"Accuracy: {acc:.4f}")
    print(f"Precision: {precision:.4f}")
    print(f"Recall: {recall:.4f}")
    print(f"F1-Score: {f1:.4f}")
    print(f"AUC-ROC: {auc:.4f}")
    print(f"Confusion Matrix:\n{cm}")
    
    return {
        'Model': model_name,
        'Accuracy': acc,
        'Precision': precision,
        'Recall': recall,
        'F1': f1,
        'AUC': auc
    }

# Evaluate all models
results = []
results.append(evaluate_model(
    y_test, y_pred_dt, best_dt.predict_proba(X_test), 'Decision Tree'
))
results.append(evaluate_model(
    y_test, y_pred_nb, nb.predict_proba(X_test), 'Naive Bayes'
))
# ... (similar for other models)
\end{verbatim}

\subsection{10-Fold Cross-Validation}

\begin{verbatim}
from sklearn.model_selection import cross_val_score

def cross_validate_model(model, X, y, model_name):
    cv_scores = cross_val_score(
        model, X, y, cv=10, scoring='accuracy', n_jobs=-1
    )
    print(f"\n{model_name} 10-Fold CV:")
    print(f"Mean Accuracy: {cv_scores.mean():.4f} ± {cv_scores.std():.4f}")
    print(f"Fold Scores: {cv_scores}")
    return cv_scores.mean(), cv_scores.std()

# Cross-validate all models
cross_validate_model(best_dt, X_scaled, y, 'Decision Tree')
cross_validate_model(nb, X_scaled, y, 'Naive Bayes')
# ... (similar for other models)
\end{verbatim}

\section{Explainable AI Implementation}

\subsection{SHAP Analysis}

\begin{verbatim}
import shap

# Random Forest SHAP
explainer_rf = shap.TreeExplainer(rf)
shap_values_rf = explainer_rf.shap_values(X_test)

# Summary plot
shap.summary_plot(
    shap_values_rf[0],  # Class 0 (Dropout)
    X_test, feature_names=X.columns
)

# Feature importance plot
shap.summary_plot(
    shap_values_rf[0], X_test, 
    feature_names=X.columns, plot_type='bar'
)

# Individual prediction explanation
shap.force_plot(
    explainer_rf.expected_value[0],
    shap_values_rf[0][0],
    X_test.iloc[0],
    feature_names=X.columns
)
\end{verbatim}

\subsection{Visualization Generation}

\begin{verbatim}
import matplotlib.pyplot as plt
import seaborn as sns

# Confusion matrix heatmap
plt.figure(figsize=(10, 8))
sns.heatmap(cm, annot=True, fmt='d', cmap='Blues')
plt.title('Confusion Matrix - Random Forest')
plt.ylabel('Actual')
plt.xlabel('Predicted')
plt.savefig('outputs/figures/confusion_matrix_rf.png', dpi=300)

# ROC curve
from sklearn.preprocessing import label_binarize
from sklearn.metrics import roc_curve, auc

y_test_bin = label_binarize(y_test, classes=[0, 1, 2])
y_score = rf.predict_proba(X_test)

plt.figure(figsize=(10, 8))
for i in range(3):
    fpr, tpr, _ = roc_curve(y_test_bin[:, i], y_score[:, i])
    roc_auc = auc(fpr, tpr)
    plt.plot(fpr, tpr, label=f'Class {i} (AUC = {roc_auc:.2f})')

plt.plot([0, 1], [0, 1], 'k--', label='Random (AUC = 0.50)')
plt.xlabel('False Positive Rate')
plt.ylabel('True Positive Rate')
plt.title('ROC Curves - Multi-Class')
plt.legend()
plt.savefig('outputs/figures/roc_curves.png', dpi=300)
\end{verbatim}

\section{Challenges and Solutions}

\subsection{Class Imbalance}

\textbf{Challenge}: Enrolled class underrepresented (794 vs. 2,209 Graduate).

\textbf{Solution}: 
\begin{itemize}
\item Stratified sampling in train-test split
\item Class weights in loss function: $w_c = \frac{N}{3 \cdot N_c}$
\item SMOTE oversampling considered but not used (degraded precision)
\end{itemize}

\subsection{Feature Correlation}

\textbf{Challenge}: Some features highly correlated (e.g., 1st sem approved vs. 1st sem grade: r=0.89).

\textbf{Solution}:
\begin{itemize}
\item Correlation analysis before training
\item Retain both features (ensemble methods handle multicollinearity)
\item Monitor VIF (Variance Inflation Factor) $<$ 10
\end{itemize}

\subsection{Computational Resources}

\textbf{Challenge}: SHAP calculation for deep learning models computationally expensive.

\textbf{Solution}:
\begin{itemize}
\item Use 100 background samples (instead of full training set)
\item GPU acceleration for neural network SHAP
\item Batch processing for large test sets
\end{itemize}

\section{Summary}

This chapter detailed the practical implementation covering data preprocessing (missing value handling, encoding, scaling), feature ranking (5 methods), baseline model training (6 algorithms with hyperparameter tuning), evaluation metrics computation (accuracy, precision, recall, F1, AUC-ROC, 10-fold CV), and explainable AI integration (SHAP analysis with visualizations). All code follows best practices with reproducible random seeds and documented configurations.
